\scanpage{36}
% Continuation.
gr.duées anticommutatives, $f^*$ homogène. Supposons qu'il existe un entier $d=d(A)$ tel que $\unsure{k|\kappa}_A$ soit composé
\diagram{\begin{tikzcd}
A\arrow[r,dash]&A^d\arrow[r,"\unsure{k|\kappa}_A^{(d)}"]&k
\end{tikzcd}}
et de même un $d'=d(B)$. Alors si $x$ est de degré $i$, la forme $\varphi_A(x)$ est nulle en degré $\ne d-i$, donc
\[
\langle f^*(y),\varphi_A(x)\rangle=\unsure{k|\kappa}(f^*(y)x)
\]
est nul en degré $\ne\unsure{d-i'|d-i}$;
\deleted{donc si $f$} de $=$, \inserted{si $x'$ \deleted{$\unsure{{}=f_*(x)}$} est hom de degré $i'$}, la forme $\varphi_B(x')$ est nulle en degré $\ne d'-i'$, donc on doit avoir (si $f_*(x)\ne0$)
\[
d-i=d'-i'\qquad\text{i.e.}\qquad i'=i+(d'-d).
\]
Si alors on pose $f_*(x)=\sum_jx'_j$, on voit que l'effet de la forme $\varphi_A(f_*(x))$ sur $B^{d'-j}$ est égal à $\varphi_B(x'_j)$, donc est nul si $\deleted{j}\inserted{d'-j}\ne d-i$, donc dans ce cas $x'_j=0$. Donc $f_*$ \uline{augmente les degrés de} $\unsure{d'-d|d'+d}$.

Nous travaillons dans la catégorie \inserted{bivariante} des Algèbres graduées anticommutatives, munies d'un degré $d$, le $f_*:A\longrightarrow B$ étant assujetti à être de degré $d(B)-d(A)$, \inserted{et $\unsure{k|\kappa}_A$ devant factoriser par $A^{d(A)}$}. La sous catégorie pleine des $A$ tels que $d(A)$ pair est commode [et il y a lieu parfois de se borner aux Algèbres $A$ qui sont à degrés \deleted{dans}\inserted{dans} $[0,d(A)]$]. Il n'est pas \deleted{utile}\inserted{commode} de se borner aux ``algèbres de Poincaré'', \unsure{associée|associées} comme dessus à une forme sur $A$, car les constructions universelles qu'on va faire en feront de tte façon sortir!
\mspage{35}
